% =========================================================================
% Capítulo 5: Mecánica Lagrangiana
% Tema 01: Mecánica Lagrangiana
% Problema 07: Dos masas en varillas paralelas unidas por un resorte
% =========================================================================
\subsubsection{Dos masas en varillas paralelas unidas por un resorte}

Dos masas idénticas $m$, etiquetadas $P$ y $Q$, están ensartadas en dos varillas
horizontales fijas y paralelas, separadas por una distancia vertical $\ell_0$, como se
muestra en la Figura \ref{fig:c5_e01_p07_masas_paralelas}. Las masas deslizan sin
fricción a lo largo de las varillas. Un resorte ideal de constante elástica $K$ y
longitud natural nula conecta ambas masas. El sistema está bajo un campo gravitatorio
uniforme $\mathbf{g}$ dirigido hacia abajo. Se definen las coordenadas generalizadas
$x_1$ y $x_2$ como las posiciones horizontales de las masas $Q$ y $P$, respectivamente,
medidas desde el extremo izquierdo de las varillas.

\begin{enumerate}
    \item Escriba las posiciones y velocidades de ambas masas.
    \item Construya la función Lagrangiana $\mathcal{L}=T-V$ del sistema.
    \item Obtenga las ecuaciones de movimiento para $x_1$ y $x_2$.
\end{enumerate}

\begin{figure}[!ht]
    \centering
    \begin{tikzpicture}[>=Latex, scale=1.2]

        % --- PAREDES ---
        \fill[pattern=north east lines] (-0.3,-1) rectangle (0,3);
        \draw[thick] (0,-1) -- (0,3);
        \fill[pattern=north east lines] (6,-1) rectangle (6.3,3);
        \draw[thick] (6,-1) -- (6,3);

        % --- VARILLAS ---
        \shade[top color=gray!50, bottom color=gray!50, middle color=gray!10] (0,-0.1) rectangle (6,0.1);
        \draw (0,-0.1) rectangle (6,0.1);
        \shade[top color=gray!50, bottom color=gray!50, middle color=gray!10] (0,1.9) rectangle (6,2.1);
        \draw (0,1.9) rectangle (6,2.1);

        % --- DISTANCIA VERTICAL l_0 ---
        \draw[<->, thick] (-0.6,0) -- (-0.6,2) node[midway, fill=white, inner sep=2pt] {$\ell_0$};

        % --- COORDENADAS x1, x2 ---
        \draw[|->, thick] (0,-0.8) -- (2,-0.8) node[midway, below] {$x_1$};
        \draw[|->, thick] (0,2.8) -- (4,2.8) node[midway, above] {$x_2$};

        % --- RESORTE ---
        \draw[thick, decoration={coil, aspect=0.5, segment length=7pt, amplitude=6pt, pre length=0.4cm, post length=0.4cm}, decorate] (2,0) -- (4,2);
        \node at (3.3, 0.8) {$K$};

        % --- MASAS ---
        \draw[fill=gray!20, rounded corners=4pt] (1.4,-0.25) rectangle (2.6,0.25);
        \fill (2,0) circle (2pt) node[above left=1pt] {$Q$};
        \node[below] at (2,-0.25) {$m$};

        \draw[fill=gray!20, rounded corners=4pt] (3.4,1.75) rectangle (4.6,2.25);
        \fill (4,2) circle (2pt) node[below right=1pt] {$P$};
        \node[above] at (4,2.25) {$m$};

        % --- GRAVEDAD ---
        \draw[->, thick] (7,1.5) -- (7,0.5) node[midway, left] {$\mathbf{g}$};

    \end{tikzpicture}
    \caption{Masas $P$ y $Q$ en varillas paralelas separadas $\ell_0$, unidas por un resorte de longitud natural nula.}
    \label{fig:c5_e01_p07_masas_paralelas}
\end{figure}

% --- SOLUCIÓN ---------------------------------------------------------------
\begin{solucion}{Dos masas en varillas paralelas unidas por un resorte}
\subsubsection*{Pregunta 1: Posiciones y velocidades}

Con el extremo izquierdo de la varilla inferior como origen $(0,0)$, las posiciones de las
masas $Q$ y $P$ son
$$\vec{r}_Q = x_1\,\hat{\imath},\qquad \vec{r}_P = x_2\,\hat{\imath} + \ell_0\,\hat{\jmath},$$
y derivando respecto al tiempo,
$$\vec{v}_Q = \dot{x}_1\,\hat{\imath},\qquad \vec{v}_P = \dot{x}_2\,\hat{\imath}.$$

\subsubsection*{Pregunta 2: Lagrangiano}

La energía cinética total es la suma de las traslaciones de ambas masas:
$$T = \tfrac12 m\dot{x}_1^2 + \tfrac12 m\dot{x}_2^2.$$

La energía potencial combina el término gravitatorio (constante, pues $P$ permanece a
altura $\ell_0$) y el elástico. Como el resorte tiene longitud natural nula, su energía es
$\tfrac12 K d^2$ con $d^2=(x_2-x_1)^2+\ell_0^2$ (teorema de Pitágoras):
$$V = m g\ell_0 + \tfrac12 K\big[(x_2-x_1)^2 + \ell_0^2\big].$$

El Lagrangiano $\mathcal{L}=T-V$ resulta
$$\boxed{\;\mathcal{L} = \tfrac12 m(\dot{x}_1^2 + \dot{x}_2^2)
- m g\ell_0 - \tfrac12 K\big[(x_2-x_1)^2 + \ell_0^2\big]\;}$$
Los términos $m g\ell_0$ y $\tfrac12 K\ell_0^2$ son constantes y no contribuyen a las
ecuaciones de movimiento.

\subsubsection*{Pregunta 3: Ecuaciones de movimiento}

Aplicando Euler-Lagrange a cada coordenada:
$$m\ddot{x}_1 = K(x_2-x_1),\qquad m\ddot{x}_2 = -K(x_2-x_1),$$
es decir,
$$\boxed{\begin{cases} m\ddot{x}_1 + K(x_1-x_2) = 0 \\[2pt] m\ddot{x}_2 + K(x_2-x_1) = 0\end{cases}}$$

Sumando ambas se obtiene $m(\ddot{x}_1+\ddot{x}_2)=0$: el centro de masa se mueve con
velocidad constante. Restándolas, la coordenada relativa $u=x_2-x_1$ satisface
$\ddot{u} + \tfrac{2K}{m}u = 0$, un movimiento armónico simple de frecuencia
$\omega=\sqrt{2K/m}$.
\end{solucion}
